\chapter{水蒸气及其动力循环}

\begin{problem}
水蒸气的定压形成过程包括哪几个过程？其\(T\)-\(s\)图上过程线下的面积表示了什么？在其汽化阶段，工质温度、焓熵等是否变化？
\end{problem}

\begin{solution}
\begin{enumerate}
  \item \textbf{形成过程}：水蒸气的定压形成过程包括预热阶段、汽化阶段和过热阶段。
  \item \textbf{面积含义}：在\(T\)-\(s\)图上，过程线下的面积代表了 \qty{1}{kg} 工质在过程中所吸收（或放出）的热量。具体来说，预热过程线下的面积代表预热热，汽化过程线下的面积表示汽化潜热，过热过程线下的面积表示过热热。
  \item \textbf{汽化阶段参数变化}：在汽化阶段，工质的温度保持为对应的饱和温度不变。但是随着水分不断变为蒸汽，其对应的比容、焓和熵均会逐渐增加。
\end{enumerate}
\end{solution}

\begin{problem}
何谓水的临界点？其临界压力、温度分别是多少？当水的压力等于或超过临界压力时，水蒸气的定压形成过程有什么不同？
\end{problem}

\begin{solution}
\begin{enumerate}
  \item \textbf{临界点}：水的临界点（C点）是饱和水线与干饱和蒸汽线的交汇点。
  \item \textbf{临界参数}：其临界温度为 \qty{374.15}{\degreeCelsius}，临界压力为 \qty{22.129}{\mega\pascal}。
  \item \textbf{形成过程的区别}：当压力等于或超过临界压力时，饱和水线与干饱和蒸汽线不再分离，汽化潜热为零。此时定压加热过程不再经历明显的汽化阶段（无湿蒸汽区），液态水在加热中直接连续地转变为过热蒸汽。
\end{enumerate}
\end{solution}

\begin{problem}
工质在锅炉中的吸热量和在汽轮机中的做功量如何计算？
\end{problem}

\begin{solution}
\begin{enumerate}
  \item \textbf{锅炉吸热量}：锅炉属于换热器，工质在其中进行定压流动吸热，吸入的热量等于其焓的增加。计算公式为：\(q=h_{2}-h_{1}\)（在标准朗肯循环中常表示为 \(q_{1}=h_{1}-h_{4}\)）。
  \item \textbf{汽轮机做功量}：工质在绝热情况下流经汽轮机时对外所做的功等于工质的焓降。计算公式为：\(w_{t}=h_{1}-h_{2}\)。
\end{enumerate}
\end{solution}

\begin{problem}
喷管的作用是什么？如何根据压力比的大小选择喷管？为什么？
\end{problem}

\begin{solution}
\begin{enumerate}
  \item \textbf{作用}：喷管是一种使流体速度得以提高的热力设备。
  \item \textbf{选择与原因}：
  \begin{enumerate}
    \item 当压力比 \(p_{2}/p_{1}\ge\beta_{c}\) 时，应采用渐缩喷管。因为在此压比下，渐缩喷管出口能获得亚音速或音速流动。
    \item 当压力比 \(p_{2}/p_{1}<\beta_{c}\) 时，应采用渐缩渐扩（拉伐尔）喷管。因为在此压比下，该喷管出口能获得超音速流动，且在最小界面处的流速理论上等于当地音速。（注：\(\beta_{c}\) 为临界压力比）。
  \end{enumerate}
\end{enumerate}
\end{solution}

\begin{problem}
何谓绝热节流？工质经过绝热节流后的效果是什么？
\end{problem}

\begin{solution}
\begin{enumerate}
  \item \textbf{定义}：流体在管道中流动时，如果流经阀门、挡板、孔板等障碍物，流体则产生涡流和摩擦，即产生局部阻力损失，因而引起压力显著下降的现象，称为节流。因节流进行得很快，工质来不及对外散热，故一般认为是绝热节流。
  \item \textbf{效果}：绝热节流前后工质的焓保持相等（\(h_{2}=h_{1}\)）。过程具有不可逆性，核心效果是引起压力下降（\(p_{2}<p_{1}\)），熵增加（\(s_{2}>s_{1}\)），最终导致工质的作功能力下降。对实际气体还会产生温度变化（热效应、零效应或冷效应）。
\end{enumerate}
\end{solution}

\begin{problem}
何谓蒸汽动力循环？为什么蒸汽动力循环不采用卡诺循环而采用朗肯循环？朗肯循环是由哪几个过程组成的？它们分别在什么设备中完成？请在\(T\)-\(s\)图上熟练地画出朗肯循环。
\end{problem}

\begin{solution}
\begin{enumerate}
  \item \textbf{概念}：动力循环是指工质连续不断地将从高温热源取得的热量的一部分转换成对外净功的过程。
  \item \textbf{采用朗肯循环的原因}：蒸汽动力循环不采用卡诺循环是因为其在技术上难以实现且存在缺陷：理论上的循环效率不高；膨胀终点的湿度太大，对汽轮机的安全不利；需要的压缩机很大且两相压缩在技术上难以实现。因此采用简单蒸汽动力装置的理想可逆循环，即朗肯循环。
  \item \textbf{过程与设备}：
  \begin{enumerate}
    \item \(1\)—\(2\) 过程：过热蒸汽的绝热膨胀作功过程，在\textbf{汽轮机}内完成。
    \item \(2\)—\(3\) 过程：乏汽的定压放热过程，在\textbf{凝汽器}中完成。
    \item \(3\)—\(4\) 过程：凝结水的绝热压缩过程，在\textbf{给水泵}中完成。
    \item \(4\)—\(5\)—\(6\)—\(1\) 过程：给水的定压吸热过程，在\textbf{锅炉}（及省煤器、过热器）中完成。
  \end{enumerate}
\end{enumerate}

（注：受纯文本环境限制，无法直接绘制图形。标准的朗肯循环\(T\)-\(s\)图请参考幻灯片第32页）
\end{solution}

\begin{problem}
中间再热循环及回热循环与简单朗肯循环相比较，有哪些区别？再热和回热的主要作用分别是什么？
\end{problem}

\begin{solution}
\begin{enumerate}
  \item \textbf{再热循环}：
  \begin{enumerate}
    \item \textbf{区别}：新蒸汽在汽轮机内作了一部分功以后，从汽轮机抽出，通过管道送到锅炉再热器中再一次进行加热，然后再送回汽轮机的中、低压缸中继续作功。
    \item \textbf{主要作用}：降低了汽轮机排汽的湿度，给提高初压创造了条件；若选取再热压力合适，还可以提高循环热效率。
  \end{enumerate}
  \item \textbf{回热循环}：
  \begin{enumerate}
    \item \textbf{区别}：利用从汽轮机的某些中间部位抽出部分做过功的蒸汽，直接送入回热加热器中来加热锅炉给水，而不是全部排入凝汽器。
    \item \textbf{主要作用}：提高了锅炉中水的预热起点温度，使得工质在锅炉内的平均吸热温度也得到提高，从而使整个循环的热效率得以提升。
  \end{enumerate}
\end{enumerate}
\end{solution}

\section*{习题}

\begin{problem}
泵只能对液态水进行加压。如果锅炉给水温度为 \qty{264}{\degreeCelsius}，水泵入口水的压力至少要大于多少兆帕？当给水泵出口压力为 \qty{17}{\mega\pascal} 时，水处于什么状态？水在 \qty{17}{\mega\pascal} 定压力下，送入锅炉受热面定压加热，当水温上升 \qty{70}{\degreeCelsius} 时，水处于什么状态？当水中产生 \(20\%\) 蒸汽时，处于什么状态？温度是多少？继续定压加热到什么状态时温度才开始上升？当温度达到 \qty{545}{\degreeCelsius} 时，处于什么状态？具有多大的过热度？
\end{problem}

\begin{solution}
（参考 PPT \textbf{【例 2.1】}）

\begin{itemize}
  \item \textbf{入口压力}：查附表可知，当饱和温度 \(t_s = \qty{264}{\degreeCelsius}\) 时，给水泵入口压力至少应大于此温度对应的饱和压力 \(p_s\)。经查表换算，压力至少要大于 \qty{5.05}{\mega\pascal}。
  \item \textbf{给水泵出口状态}：当压力为 \qty{17}{\mega\pascal} 时，其对应的饱和温度 \(t_s \approx \qty{352.37}{\degreeCelsius}\)。因为实际给水温度 \qty{264}{\degreeCelsius} \(< t_s\)，故水处于\textbf{过冷水状态}。
  \item \textbf{水温上升 70°C}：加热后温度为 \qty{334}{\degreeCelsius}。因为 \qty{334}{\degreeCelsius} \(< t_s\)，此时水仍处于\textbf{过冷水状态}。
  \item \textbf{产生 20\% 蒸汽}：当给水中产生蒸汽时（干度 \(x = 0.2\)），水处于\textbf{湿饱和蒸汽状态}。此时温度恒等于饱和温度，即 \qty{352.37}{\degreeCelsius}。
  \item \textbf{温度开始上升的节点}：继续定压加热，当液态水完全汽化，达到\textbf{干饱和蒸汽状态}（即工质焓值达到相应压力下的干饱和蒸汽焓 \(h''\)，约为 \qty{2532.6}{\kilo\joule\per\kilo\gram}）时，其温度才开始升高。
  \item \textbf{达到 545°C 的状态与过热度}：此时工质温度超过饱和温度，处于\textbf{过热蒸汽状态}。其过热度为
  \begin{equation*}
    D = t - t_s = 545 - 352.37 = \qty{192.63}{\degreeCelsius}
  \end{equation*}
\end{itemize}
\end{solution}

\begin{problem}
某电站锅炉参数如下，高压锅炉是：锅炉出口过热蒸汽压力 \(p_1 = \qty{10}{\mega\pascal}\)，温度 \(t_1 = \qty{540}{\degreeCelsius}\)，给水温度 \(t_2 = \qty{220}{\degreeCelsius}\)；亚临界压力锅炉是：\(p_1 = \qty{17}{\mega\pascal}\)，\(t_1 = \qty{540}{\degreeCelsius}\)，\(t_2 = \qty{260}{\degreeCelsius}\)。按定压加热，分别将炉内加热过程定性地描绘在 \(T\)-\(s\) 图上，并比较其预热热、汽化潜热和过热热因 \(p_1\) 不同所造成的变化趋势。
\end{problem}

\begin{solution}
（参考 PPT 中 \textbf{\(T\)-\(s\) 图}及定压加热过程规律）

在 \(T\)-\(s\) 图上，这两个加热过程均沿定压线向右上侧延伸，经过预冷水、湿饱和蒸汽区，最终到达同为 \qty{540}{\degreeCelsius} 的过热区。当 \(p_1\) 从 \qty{10}{\mega\pascal} 提高到 \qty{17}{\mega\pascal} 时，各部分热量变化趋势如下：

\begin{itemize}
  \item \textbf{预热热（\(q'\)）增加}：随着压力升高，饱和水线（MC）整体向右上方偏移，饱和温度 \(t_s\) 随之升高，将水从过冷态加热至饱和水状态所需的预热热线下面积显著增大。
  \item \textbf{汽化潜热（\(r\)）减小}：在 \(T\)-\(s\) 图中，随着压力升高并逼近临界点，饱和水线与干饱和蒸汽线（NC）之间的水平距离收缩，导致汽化过程线下的面积（潜热）减小。
  \item \textbf{过热热（\(q''\)）减小}：由于两者的终点过热温度相同（均为 \qty{540}{\degreeCelsius}），但亚临界锅炉的饱和温度 \(t_s\) 更高，导致将蒸汽加热到规定温度所需的过热温差缩小。因此，过热过程线下的面积减小。
\end{itemize}
\end{solution}

\begin{problem}
如在汽轮机的某级中，蒸汽以 \(p_1 = \qty{0.5}{\mega\pascal}\)，\(t_1 = \qty{230}{\degreeCelsius}\) 进入喷管，喷管的背压 \(p_b = \qty{0.4}{\mega\pascal}\)，蒸汽流量为 \qty{25}{\tonne\per\hour}，试问应选择何种型式的喷管？并计算喷管出口的流速 \(c_2\) 和喷管出口截面积 \(S_2\)。
\end{problem}

\begin{solution}
（参考 PPT \textbf{【例 2.3】}与\textbf{通过喷管的绝热流动}）

\begin{itemize}
  \item \textbf{喷管型式选择}：已知初始状态为过热蒸汽，过热蒸汽的临界压力比 \(\beta_c = 0.546\)。实际的压力比为 \(p_2/p_1 = 0.4 / 0.5 = 0.8\)。因为 \(0.8 > 0.546\)，喷管出口只能获得亚音速，故必须选用\textbf{渐缩形喷管}。

  \item \textbf{出口流速 \(c_2\)}：忽略初速度，根据绝热流动能量方程，流速计算公式为 \(c_2 = 44.72\sqrt{h_1 - h_2}\)。

  查表得 \(p_1 = \qty{0.5}{\mega\pascal}\)，\(t_1 = \qty{230}{\degreeCelsius}\) 时，初焓 \(h_1 \approx \qty{2918}{\kilo\joule\per\kilo\gram}\)，熵 \(s_1 \approx \qty{7.059}{\kilo\joule\per\kilo\gram\per\kelvin}\)。

  沿等熵过程绝热膨胀至 \(p_2 = \qty{0.4}{\mega\pascal}\)，查表得出口焓 \(h_2 \approx \qty{2820}{\kilo\joule\per\kilo\gram}\)，比容 \(v_2 \approx \qty{0.52}{\cubic\meter\per\kilo\gram}\)。

  代入公式验算得：
  \begin{equation*}
    c_2 = 44.72\sqrt{2918 - 2820} \approx \qty{442}{\meter\per\second}
  \end{equation*}

  \item \textbf{出口截面积 \(S_2\)}：蒸汽的质量流量 \(q_m = 25 \times 10^3 / 3600 \approx \qty{6.94}{\kilo\gram\per\second}\)。

  截面积公式计算：
  \begin{equation*}
    S_2 = \frac{q_m \cdot v_2}{c_2} = \frac{6.94 \times 0.52}{442} \approx \qty{0.00817}{\square\meter}
  \end{equation*}
  （即 \qty{81.7}{\square\centi\meter}）。
\end{itemize}
\end{solution}

\begin{problem}
若基本朗肯循环的初参数 \(p_1 = \qty{13.5}{\mega\pascal}\)，\(t_1 = \qty{550}{\degreeCelsius}\)，凝汽器压力 \(p_2 = \qty{5}{\kilo\pascal}\)，不计水泵功。试求：
\begin{enumerate}
  \item 该循环的 \(w_0\)、\(q_1\)、\(x\)、\(\eta_t\)；
  \item 如果工质在汽轮机内膨胀到 \qty{3.5}{\mega\pascal} 时，再去加热到 \qty{550}{\degreeCelsius}，此再热循环的 \(w_0\)、\(q_1\)、\(x\)、\(\eta_{TR}\)；
  \item 比较两种循环方式下的热效率和乏汽干度。
\end{enumerate}
\end{problem}

\begin{solution}
（参考 PPT \textbf{【例 2.4】}与\textbf{【例 2.5】}）

\begin{enumerate}
  \item \textbf{基本朗肯循环}：
  \begin{itemize}
    \item \textbf{计算方法}：循环作功 \(w_0 = h_1 - h_2\)，吸热量 \(q_1 = h_1 - h_3\)（忽略水泵功），热效率 \(\eta_T = \frac{h_1 - h_2}{h_1 - h_3}\)。
    \item \textbf{参数验算}：查表得初态 \(h_1 \approx \qty{3447}{\kilo\joule\per\kilo\gram}\)，\(s_1 \approx \qty{6.54}{\kilo\joule\per\kilo\gram\per\kelvin}\)。等熵膨胀至 \qty{5}{\kilo\pascal}，查得排汽干度 \(x \approx 0.766\)，排汽焓 \(h_2 \approx \qty{1994}{\kilo\joule\per\kilo\gram}\)，饱和水焓 \(h_3 \approx \qty{138}{\kilo\joule\per\kilo\gram}\)。
    \item \textbf{结果}：
    \begin{equation*}
      w_0 = 3447 - 1994 = \qty{1453}{\kilo\joule\per\kilo\gram}
    \end{equation*}
    \begin{equation*}
      q_1 = 3447 - 138 = \qty{3309}{\kilo\joule\per\kilo\gram}
    \end{equation*}
    \begin{equation*}
      \eta_t = \frac{1453}{3309} \approx 43.9\%
    \end{equation*}
  \end{itemize}

  \item \textbf{再热循环}：
  \begin{itemize}
    \item \textbf{计算方法}：从高压缸排汽点 \(b\) 抽汽再热至点 \(a\)。净功 \(w_{\text{net}} = (h_1 - h_b) + (h_a - h_2)\)，吸热量 \(q_1 = (h_1 - h_4) + (h_a - h_b)\)，再热循环热效率 \(\eta_{t,RH} = \frac{w_{\text{net}}}{q_1}\)。
    \item \textbf{参数验算}：工质等熵膨胀至 \qty{3.5}{\mega\pascal} 时，查得 \(h_b \approx \qty{3020}{\kilo\joule\per\kilo\gram}\)。再热至 \qty{550}{\degreeCelsius} 后，\(h_a \approx \qty{3568}{\kilo\joule\per\kilo\gram}\)。等熵膨胀至 \qty{5}{\kilo\pascal}，新排汽干度 \(x \approx 0.86\)，排汽焓 \(h_2 \approx \qty{2221}{\kilo\joule\per\kilo\gram}\)。
    \item \textbf{结果}：
    \begin{equation*}
      w_{\text{net}} = (3447 - 3020) + (3568 - 2221) = \qty{1774}{\kilo\joule\per\kilo\gram}
    \end{equation*}
    \begin{equation*}
      q_1 = (3447 - 138) + (3568 - 3020) = \qty{3857}{\kilo\joule\per\kilo\gram}
    \end{equation*}
    \begin{equation*}
      \eta_{TR} = \frac{1774}{3857} \approx 46.0\%
    \end{equation*}
  \end{itemize}

  \item \textbf{比较结论}：
  \begin{itemize}
    \item 采用一次中间再热后，\textbf{循环热效率得到了提高}（\(46.0\% > 43.9\%\)）。
    \item \textbf{乏汽干度显著改善}（从 \(0.766\) 提升至 \(0.86\)），这不仅为进一步提高初压创造了条件，也降低了末级叶片的水滴冲蚀，对汽轮机的安全运行极为有利。
  \end{itemize}
\end{enumerate}
\end{solution}

\begin{problem}
设有两级抽汽回热循环，其参数 \(p_1 = \qty{10}{\mega\pascal}\)，\(t_1 = \qty{540}{\degreeCelsius}\)，凝汽器压力 \(p_2 = \qty{5}{\kilo\pascal}\)，汽轮机第一级、第二级抽汽压力分别为 \(p_{01} = \qty{0.4}{\mega\pascal}\) 和 \(p_{02} = \qty{0.21}{\mega\pascal}\)。试求该两级回热循环的汽轮机内功和循环热效率。
\end{problem}

\begin{solution}
（参考 PPT \textbf{两级抽汽回热循环的热效率}）

\begin{itemize}
  \item \textbf{建立回热器热量与质量平衡}：假设进入汽轮机的新蒸汽为 \qty{1}{\kilo\gram}，分别抽出质量比例为 \(\alpha_1\) 和 \(\alpha_2\) 的蒸汽送入 I 级和 II 级回热加热器中。

  根据混合式回热器的热平衡方程：
  \begin{itemize}
    \item I 级回热器：\(\alpha_1(h_{01} - h_{02}') = h_{01}' - h_{02}'\)。
    \item II 级回热器：\(\alpha_2(h_{02} - h_2') = (1 - \alpha_1)(h_{02}' - h_2')\)。
  \end{itemize}

  \item \textbf{状态参数获取与抽汽比验算}：

  初态 \(h_1 \approx \qty{3475}{\kilo\joule\per\kilo\gram}\)；各级抽汽焓 \(h_{01} \approx \qty{2690}{\kilo\joule\per\kilo\gram}\)，\(h_{02} \approx \qty{2550}{\kilo\joule\per\kilo\gram}\)；排汽焓 \(h_2 \approx \qty{2030}{\kilo\joule\per\kilo\gram}\)。对应的各级饱和水焓约为 \(h_{01}' \approx \qty{605}{\kilo\joule\per\kilo\gram}\)，\(h_{02}' \approx \qty{510}{\kilo\joule\per\kilo\gram}\)，\(h_2' \approx \qty{138}{\kilo\joule\per\kilo\gram}\)。

  解得抽汽系数：\(\alpha_1 \approx 0.044\)；\(\alpha_2 \approx 0.147\)。剩余排入凝汽器的份额 \(\alpha_0 = 1 - 0.044 - 0.147 = 0.809\)。

  \item \textbf{汽轮机内功（\(w_0\)）计算}：
  \begin{equation*}
    w_0 = \alpha_1(h_1 - h_{01}) + \alpha_2(h_1 - h_{02}) + \alpha_0(h_1 - h_2)
  \end{equation*}

  代入数值验算：
  \begin{equation*}
    w_0 = 0.044 \times 785 + 0.147 \times 925 + 0.809 \times 1445 \approx \qty{1339.5}{\kilo\joule\per\kilo\gram}
  \end{equation*}

  \item \textbf{循环热效率（\(\eta_{TRG}\)）计算}：

  循环总加热量：
  \begin{equation*}
    q_1 = h_1 - h_{01}' = 3475 - 605 = \qty{2870}{\kilo\joule\per\kilo\gram}
  \end{equation*}

  热效率公式：
  \begin{equation*}
    \eta_{TRG} = 1 - \frac{\alpha_0(h_2 - h_2')}{h_1 - h_{01}'}
  \end{equation*}

  代入数值验算：
  \begin{equation*}
    \eta_{TRG} = \frac{1339.5}{2870} \approx 46.7\%
  \end{equation*}

  \item \textbf{结论}：被抽出的 \(\alpha_1\)、\(\alpha_2\) 部分蒸汽作了功而没有向外放热（相当于该部分热效率达到了100\%），最终使锅炉中水的预热起点温度提高，工质在锅炉内的平均吸热温度也大幅提高，进而提升了整个循环的热效率。
\end{itemize}
\end{solution}
